% Responsable: @victoriaguilllen · Revisora: @gvidalquadras (issue #18)
% Extensión: 1–1,5 págs. Figuras: exp. 1 (gap vs n) y exp. 2 (stable rank vs n).
\section{Gap espectral y colapso de rango en anchura}\label{sec:gap}

\pendiente{Definiciones y resultado principal.}

% Contenido previsto:
%  - Definición de matriz estocástica por filas; por qué sigma_1(A) = 1 (dirección \ones).
%  - Modelo aleatorio en inicialización (pesos N(0, sigma_qk^2), X X^T = I).
%  - Teorema: sigma_1(A) -> 1, sqrt(n) sigma_2(A) -> 2 sigma.
%  - Definición de stable rank: \sr(M) = \normF{M}^2 / \norm{M}^2.
%  - Intuición de la prueba y por qué A ≈ (1/n) \ones\ones^T implica colapso en anchura.
%
% \begin{figure}[t]
%   \centering
%   \begin{subfigure}{0.48\linewidth}
%     \includegraphics[width=\linewidth]{exp1_gap_vs_n}
%     \caption{$\sigma_1$ y $\sigma_2$ frente a $n$.}
%   \end{subfigure}\hfill
%   \begin{subfigure}{0.48\linewidth}
%     \includegraphics[width=\linewidth]{exp2_rank_vs_n}
%     \caption{Rango estable de $AX$ frente a $n$.}
%   \end{subfigure}
%   \caption{...}\label{fig:gap}
% \end{figure}
